% page 350 du fac-simile (image p-349 du scan)
% bloc 154.9 x 233.3 mm ; marge gauche 8.0 mm
\pgc{-12.700mm}{0.000mm}{
\l{\cel{3}342}
\l{}
\l{\sou{litotriciar}. (\sou{trans.}) (\sou{kirurg.}).\cel{2}Triturar (per operaco)}
\l{\cel{3}la kalkolo qua esas en la veziko, enduktante instrumen-}
\l{\cel{3}to specala en-alonge la uretro.- DEFIRS.}
\l{}
\l{\sou{litotripsiar}. (\sou{trans}.) (\sou{kirurg.})\cel{2}Dividar per froto ed ex-}
\l{\cel{3}traktar la kalkolo ek la veziko. - DEFIRS.}
\l{}
\l{\sou{litro}. I. (en\cel{1}\sou{la mezuro-sistemo metrala}): unajo di la me-}
\l{\cel{3}zuri pri kapaceso. - II. Botelo kapaca ye un litro.}
\l{\cel{3}- DEFIRS.}
\l{}
\l{\sou{liturgio.}\cel{2}(\sou{religio)}. Formo di la kulto ; ordino di la}
\l{\cel{3}ceremonii konsakrita. - DEFIRS.}
\l{}
\l{\sou{liuto}. (\sou{Muziko)}.\cel{2}Muzik-instrumento kun plura rangi de kordi}
\l{\cel{3}muntita sur rezono-kesto sube-rondigita, e quan onu uzas}
\l{\cel{3}per pinchar la kordi. - DEFIRS.}
\l{}
\l{\sou{livar}. (\sou{trans}.) Departar de apud (lu) definitive. - E.}
\l{}
\l{\sou{livida}. (\sou{aludante la vizajo}). Qua esas plombo-nigra, blua-}
\l{\cel{3}tra, malado-pala. - DEFIS.}
\l{}
\l{\sou{livrar}. (\sou{trans., ad)}.\cel{2}I. Igar ye la dispono di ulu. -II.}
\l{\cel{3}Submisar a la ago da ulu. - EFS.}
\l{}
\l{\sou{livreo}. I. Vesto quan la reji, la siniori furnisas a la}
\l{\cel{3}viri di sua eskorto, e di qua la koloro, la galoni, la}
\l{\cel{3}butoni, e c. memorigas lia blazoni. - II. Distingo-kostumo}
\l{\cel{3}(distingebla pro koloro, galoni, butoni, e c.) quan la}
\l{\cel{3}hemestro igas +werar da sua servistuli. - DEFIRS.}
\l{}
\l{\sou{lizo}. I. Sedimento quan la vino depozas sur la fondo di}
\l{\cel{3}barelo. - II. Elemento di reziduo sociala. - EF.}
\l{}
\l{\sou{lizimako.}\cel{2}(\sou{bot.}) Planto qua formacas genero ek la familio}
\l{\cel{3}\textquotedbl{}primulacei\textquotedbl{}. - L.\cel{1}\sou{lysimachia}.}
\l{}
\l{lo. - Pronomo qua : 1. reprezentas fakto, o propoziciono ;}
\l{\cel{3}2. indikas kozo nedeterminita precize (= to quo, co}
\l{}
\l{}
\l{lobo.\cel{2}(\sou{anat}.)I. Parto di ula organi di qua la formo esas}
\l{\cel{3}ronda. - II. Saliajo ronda per qua finas, ye sua infra}
\l{\cel{3}parto, la funelo di la orelo. - EFIS.}
\l{}
\l{\sou{lobelio}. (\sou{bot.}) Genero de \textquotedbl{}lobeliacei\textquotedbl{}, di qua la suko esas}
\l{\cel{3}laktatra, akra, venena, e quan onu kultivas precipue kom}
\l{\cel{3}planto orniva. - L.\cel{1}\sou{lobelia}. - DEFS.}
\l{}
\l{\sou{locho}. (\sou{zool.}) Fisho di aquo sen-sala, di qua la korpo esas}
\l{\cel{3}tre oblonga. - L.\cel{1}\sou{cobitis barbatula}. - EFIS.}
}
